
\documentclass[../../main.tex]{subfiles}
\begin{document}

% 年度の大きな表示
\begin{center}
  {\fontsize{48pt}{58pt}\selectfont \textbf{1990}\normalsize\textbf{年度}}
\end{center}

\vspace{10mm}

% 本文

\vspace{15mm}

% 出題テーマと難易度の表
\noindent\textbf{出題テーマと難易度}

\vspace{5mm}

\begin{center}
  \shadowbox{%
    \begin{tabular}{|c|p{8cm}|c|}
      \hline
      \textbf{問題番号} & \textbf{テーマ} & \textbf{難易度} \\
      \hline
      第1問 & 不等式の等号成立条件 & \\
      \hline
      第2問 & 対数と不等式の証明 & \\
      \hline
      第3問 & 楕円と円上の動点間距離の最大値 & \\
      \hline
      第4問 & 三角関数の積分と最小値 & \\
      \hline
      第5問 & 楕円の接線と図形の面積 & \\
      \hline
    \end{tabular}%
  }
\end{center}

\vspace{15mm}

% 解答
\noindent\textbf{解答}

\vspace{5mm}

\begin{center}
  \shadowbox{%
    \renewcommand{\arraystretch}{1.6}
    \begin{tabular}{|c|c|p{8cm}|}
      \hline
      \textbf{問題番号} & \textbf{小問} & \textbf{答え} \\
      \hline
      第1問 & - & $xw=yz$ \\
      \hline
      第2問 & - & 証明問題 \\
      \hline
      第3問 & - & $6$ \\
      \hline
      第4問 & - & $\log\dfrac{2+\sqrt{2}}{2-\sqrt{2}}$ \\
      \hline
      第5問 & - & 図示，面積 $\dfrac{1}{3}\left(\dfrac{4}{3}\pi-\sqrt{3}\right)$ \\
      \hline
    \end{tabular}%
    \renewcommand{\arraystretch}{1.0}
  }
\end{center}

\vspace{15mm}

% 解答時間
\noindent\textbf{試験形態}

\vspace{5mm}

\begin{center}
  \shadowbox{%
    \begin{tabular}{p{8cm}}
      （要確認）
    \end{tabular}%
  }
\end{center}

\end{document}
